% (C) Copyright 2016
% Urs Fässler, www.bitzgi.ch
% SPDX-License-Identifier: CC-BY-SA-4.0

\subsection{Handhabung der Lizenzen}
\label{sec:handhabung}
\subsectionframe

\begin{frame}{Rechte am Code}
	\pause
	\begin{itemize}
		\item Besitzt man die Rechte?
		\pause
		\begin{itemize}
			\item Anstellungsvertrag
			\item Contributor License Agreement
		\end{itemize}
		\pause
		\item Besitzer darf mit Code machen was er will
	\end{itemize}
\end{frame}

\begin{frame}{Frei mit Proprietär: Betriebssystem}
	\only<2|handout:0>{
		\begin{center}
			\includegraphics[width=0.4\textwidth]{res/propritary-on-os-base.pdf}
		\end{center}
	}
	\only<3|handout:0>{
		\begin{center}
			\includegraphics[width=0.4\textwidth]{res/propritary-on-os-app.pdf}
		\end{center}
	}
	\only<4|handout:0>{
		\begin{center}
			\includegraphics[width=0.4\textwidth]{res/propritary-on-os-permissive.pdf}
		\end{center}
	}
	\only<5|handout:1>{
		\begin{center}
			\includegraphics[width=0.4\textwidth]{res/propritary-on-os-copyleft.pdf}
		\end{center}
	}
\end{frame}
\note
{
	\begin{itemize}
		\item proprietäre Applikation auf freiem OS (z.B. Debian GNU/Linux) kein Problem
		\item Open Handcuffs: CC-0 by qubodup \url{https://openclipart.org/detail/171929/open-handcuffs}
	\end{itemize}
}

\begin{frame}{Frei mit Proprietär: Library}
	\begin{multicols}{2}
		\only<2|handout:0>{
			\begin{center}
				\includegraphics[width=0.4\textwidth]{res/propritary-dynamic-linking-base.pdf}
			\end{center}
		}
		\only<3|handout:0>{
			\begin{center}
				\includegraphics[width=0.4\textwidth]{res/propritary-dynamic-linking-permissive.pdf}
			\end{center}
		}
		\only<4|handout:0>{
			\begin{center}
				\includegraphics[width=0.4\textwidth]{res/propritary-dynamic-linking-copyleft.pdf}
			\end{center}
		}
		\only<5-|handout:1>{
			\begin{center}
				\includegraphics[width=0.4\textwidth]{res/propritary-dynamic-linking-lgpl.pdf}
			\end{center}
		}
		\pagebreak
		\only<6|handout:0>{
			\begin{center}
				\includegraphics[width=0.4\textwidth]{res/propritary-static-linking-base.pdf}
			\end{center}
		}
		\only<7|handout:0>{
			\begin{center}
				\includegraphics[width=0.4\textwidth]{res/propritary-static-linking-permissive.pdf}
			\end{center}
		}
		\only<8|handout:0>{
			\begin{center}
				\includegraphics[width=0.4\textwidth]{res/propritary-static-linking-copyleft.pdf}
			\end{center}
		}
		\only<9-|handout:1>{
			\begin{center}
				\includegraphics[width=0.4\textwidth]{res/propritary-static-linking-lgpl.pdf}
			\end{center}
		}
	\end{multicols}
\end{frame}
\note
{
	\begin{itemize}
		\item dynamisches linken nur gegen Libraries mit Public-Domain, permissiven Lizenzen oder LGPL
		\item statisches linken nur gegen Libraries mit Public-Domain oder permissiven Lizenzen
	\end{itemize}
}

\begin{frame}{Frei mit Proprietär: Code Übernahme}
	\begin{multicols}{2}
		\begin{center}
			\only<2|handout:0>{
				\includegraphics[width=0.4\textwidth]{res/propritary-use-code-base.pdf}
			}
			\only<3|handout:0>{
				\includegraphics[width=0.4\textwidth]{res/propritary-use-code-permissive.pdf}
			}
			\only<4-|handout:1>{
				\includegraphics[width=0.4\textwidth]{res/propritary-use-code-copyleft.pdf}
			}
		\end{center}
		\pagebreak
		\only<5|handout:0>{
			\begin{center}
				\includegraphics[width=0.4\textwidth]{res/propritary-cloud-base.pdf}
			\end{center}
		}
		\only<6|handout:0>{
			\begin{center}
				\includegraphics[width=0.4\textwidth]{res/propritary-cloud-permissive.pdf}
			\end{center}
		}
		\only<7|handout:0>{
			\begin{center}
				\includegraphics[width=0.4\textwidth]{res/propritary-cloud-copyleft.pdf}
			\end{center}
		}
		\only<8-|handout:1>{
			\begin{center}
				\includegraphics[width=0.4\textwidth]{res/propritary-cloud-agpl.pdf}
			\end{center}
		}
	\end{multicols}
\end{frame}
\note
{
	\begin{itemize}
		\item Verwenden in proprietärem Code: nur Public-Domain oder permissive Lizenzen
		\item Code unter Copyleft darf nicht in proprietärer Applikation verwendet werden (ausser wenn ausschliesslich intern gebraucht, dann ist es Freie Software)
		\item Wird Software als Service zur Verfügung gestellt muss der Copyleft Quelltext nicht zur Verfügung gestellt werden. Die AGPL schliesst diese Lücke. (\url{https://de.wikipedia.org/wiki/GNU\_Affero\_General\_Public\_License})
	\end{itemize}
}

\begin{frame}{Verwendung mit Freier Software}
	\pause
	\begin{itemize}
		\item Patches unter gleicher Lizenz
		\pause
		\item Lizenz Kompatibilität beachten
	\end{itemize}
	\pause
	\begin{center}
		\includegraphics[width=\textwidth]{res/floss-license-compability.png}
	\end{center}
\end{frame}
\note
{
	\begin{itemize}
		\item Floss-license-slide-image.png: CC-SA 3.0 by David A. Wheeler (\url{https://commons.wikimedia.org/wiki/File:Floss-license-slide-image.png})
	\end{itemize}
}

\begin{frame}{Auswahl einer Lizenz}
	\begin{center}
		\begin{tabular}{ccc}
		\visible<2->{\includegraphics[width=1cm]{res/PD-icon.pdf}} & \visible<4->{\includegraphics[width=1cm]{res/by.pdf}} & \visible<6->{\includegraphics[width=1cm]{res/copyleft.pdf}} \\ 
		\visible<2->{Public Domain} & \visible<4->{Permissiv} & \visible<6->{Copyleft} \\
		\visible<3->{\includegraphics[width=2.8cm]{res/cc-zero.pdf}} & \visible<5->{\includegraphics[width=2.8cm]{res/apache-logo.pdf}} & \visible<7->{\includegraphics[width=2.8cm]{res/gpl-v3-logo.pdf}} \\
		\hspace{0.3\textwidth} & \hspace{0.3\textwidth} & \hspace{0.3\textwidth} \\
		\end{tabular} 
	\end{center}
	\begin{itemize}
		\visible<8->{\item \url{http://choosealicense.com/}}
	\end{itemize}
	\vspace{1cm}
\end{frame}
\note
{
	\begin{itemize}
		\item Was will man schützen?
		\item Was will man erlauben?
		\item Was will man verhindern?
	\end{itemize}
}

\begin{frame}{Deklaration der Lizenz}
	\begin{center}
		\only<2|handout:1>
		{
			\includegraphics[width=\textwidth]{res/lizenz-source.png}
		}
		\only<3|handout:2>
		{
			\includegraphics[width=\textwidth]{res/lizenz-file.png}
		}
		\only<4|handout:3>
		{
			\includegraphics[height=6cm]{res/lizenz-binary.png}
		}
	\end{center}
\end{frame}
\note
{
	\begin{itemize}
		\item im Projekt (COPYING)
		\item in den Files (Header, SPDX)
		\item in dem Binary (About Dialog, Info Option in CLI Applikationen)
		\item Keine Lizenz: Code darf nicht benutzt werden
	\end{itemize}
}

\begin{frame}{Veröffentlichen der Sourcen}
	\only<2|handout:1>{
		\begin{center}
			\makebox[\textwidth]{\includegraphics[width=0.98\paperwidth]{res/source-download.png}}
		\end{center}
	}
	\only<3|handout:2>{
		\begin{center}
			\makebox[\textwidth]{\includegraphics[width=0.98\paperwidth]{res/source-git.png}}
		\end{center}
	}
\end{frame}
\note
{
	\begin{itemize}
		\item Source des Programms muss in der gleichen Version zur Verfügung stehen
		\item Besser noch wenn Entwicklung offen und Zugriff auf Source Versionskontrolle besteht
	\end{itemize}
}

\begin{frame}{Verteidigung des eigenen Codes}
	\begin{itemize}
		\pause
		\item Dialog suchen
		\pause
		\item Software Freedom Conservancy \url{https://sfconservancy.org/}
	\end{itemize}
\end{frame}
\note
{
	\begin{itemize}
		\item \url{https://de.wikipedia.org/wiki/Software_Freedom_Conservancy}
		\item Erfolge vor Gericht
	\end{itemize}
}
